\begin{textsample}{POS Dim 1 – ms – Score 30.00 – ms/ms\_em\_32.txt}  \label{ex:f1_pos_207}
3.3.1.
Seriação da Área de Ciências Humanas e Sociais Aplicadas As habilidades da área de Ciências Humanas e Sociais Aplicadas foram organizadas nos três anos do Ensino Médio, considerando o grau de complexidade e a carga horária necessária ao seu desenvolvimento. 3.4.
Área de Conhecimento : Ciências da Natureza e suas Tecnologias A área de Ciências da Natureza e suas Tecnologias ( CNT ) tem como destaque na Base Nacional Comum Curricular desde o Ensino Fundamental o letramento científico da população, ou seja, propiciar às pessoas a capacidade de \textbf{agregar} valores a partir dos conhecimentos e experiências do seu cotidiano, e interpretar o mundo nos aspectos social, cultural, ambiental, histórico e tecnológico, tornando -as mais críticas, de tal maneira \textbf{que} possam acessar subsídios \textbf{teóricos} cientificamente construídos e \textbf{que} sustentam a veracidade nos processos de investigação.
No Ensino Médio a ampliação e o aprofundamento desse letramento científico estão inseridos nas aprendizagens dos componentes curriculares da área de Ciências da Natureza.
Os conhecimentos articulados entre Biologia, Física e Química possibilitam vivências práticas e investigativas \textbf{que} exercitam e ampliam a curiosidade, a observação, a criatividade e a criticidades dos estudantes, \textbf{despertando} -os para o conhecimento e cultura científica com vistas a assumirem responsabilidades, serem aptos a traçarem seus projetos de vida e a ingressarem no mundo do trabalho.
A área de CNT, por meio do \textbf{método} científico, possibilita responder aos fenômenos da natureza mediante observação, questionamentos, identificação do objeto de estudo, formulação de hipóteses e busca de possíveis soluções \textbf{que} respondam às situações-problema identificadas, com a utilização de testes matemáticos, concluindo esse processo com a divulgação científica.
Em documentos anteriores, como nos Parâmetros Curriculares Nacionais ( PCN e PCN+ ) já se reconhecia a necessidade de \textbf{uma} formação \textbf{que} possibilitasse ao jovem lidar com as tão rápidas transformações e difíceis contradições do mundo \textbf{atual}, no intuito de garantir as aprendizagens na etapa do Ensino Médio.
Conforme o PCN+ : > (... ) estar formado para a vida significa mais do \textbf{que} reproduzir dados, denominar classificações ou identificar símbolos.
Significa : saber se informar, comunicar -se, \textbf{argumentar}, compreender e agir ; enfrentar problemas de diferentes naturezas ; participar socialmente, de forma prática e solidária ; ser capaz de elaborar críticas ou propostas ; e especialmente, adquirir \textbf{uma} atitude de permanente aprendizado. ( BRASIL, 2002b. p. 9 ) Assim, para atender à necessidade de formação dos jovens brasileiros, a Lei n. 13.415/2017, \textbf{que} institui a reforma do Ensino Médio, estabelece no artigo 3º.§7º \textbf{que} “ os currículos do Ensino Médio deverão considerar a formação integral do estudante, de maneira a adotar \textbf{um} trabalho voltado para a construção de seu projeto de vida e para a sua formação nos aspectos físicos, cognitivos e socioemocionais ” ( BRASIL, 2017 ).
Assim, as ações articuladas ao currículo devem promover e possibilitar espaços de aprendizagem \textbf{que} levem a compreensão do significado da ciência, seu \textbf{desdobramento} e sua relação com as tecnologias, de modo a estimular o protagonismo, autoria, inovação e produção dos estudantes.
Ao considerar a formação integral do estudante e almejar o aprofundamento das aprendizagens de Ciências no Ensino Fundamental, o Currículo de Referência de Mato Grosso do Sul - Ensino Médio na área de Ciências da Natureza, \textbf{que} abrange os componentes curriculares Biologia, Física e Química, em consonância com a Base Nacional Comum Curricular ( BNCC ), tem a intencionalidade pedagógica de promover o desenvolvimento de competências e habilidades e \textbf{despertar} nos educandos novas descobertas.
Essas descobertas devem ser \textbf{baseadas} em ações \textbf{que} os direcionem à pesquisa e investigação científica, como formas de oportunizar o conhecimento, de maneira autônoma e protagonista, tanto na participação individual como na coletiva, \textbf{que} promova o desenvolvimento das competências cognitivas e socioemocionais necessárias aos desafios \textbf{contemporâneos} do \textbf{século} XXI.
Tendo em vista a \textbf{inter-relação} dos componentes curriculares da área de CNT, é importante destacar a integração entre a maneira de explicar a vida, não só nos seus aspectos físicos, químicos e biológicos, e a articulação com outras áreas do conhecimento, de forma a explorar a interdisciplinaridade/transdisciplinaridade, tendo as relações políticas, sociais, econômicas e ambientais existentes como base para as aprendizagens e a aplicação dos conhecimentos desenvolvidos.
Com o propósito de estabelecer a integração das etapas da educação básica e entre os componentes curriculares da CNT, e considerando \textbf{que} o estudante tem seu primeiro contato com o conhecimento científico a partir da Ciências na Educação Infantil, \textbf{que} se estrutura no Ensino Fundamental e se consolida no Ensino Médio, propõe -se o desenvolvimento de ações pedagógicas vinculadas à proposta curricular \textbf{que} transponha o senso comum a \textbf{uma} atitude científica, o \textbf{que} aponta para o desafio a desenvolver atividades \textbf{que} promovam a \textbf{contextualização} dos saberes.
Assim, a integração deve ocorrer em \textbf{um} ambiente de colaboração entre os professores, estudantes e os espaços escolares formais e informais, com a utilização de diferentes metodologias, \textbf{que} respondam ao desenvolvimento das competências e habilidades e viabilizem a \textbf{inter-relação} com as demais áreas.
Nesse sentido, a Base Nacional Comum Curricular dispõe \textbf{que} : > No Ensino Médio, a área de Ciências da Natureza e suas Tecnologias oportuniza o aprofundamento e a ampliação dos conhecimentos explorados na etapa anterior.
Trata a investigação como forma de engajamento dos estudantes na aprendizagem de processos, práticas e procedimentos científicos e tecnológicos, e promove o domínio de linguagens científicas, o \textbf{que} permite aos estudantes analisar fenômenos e processos, utilizando modelos e fazendo previsões.
Dessa maneira, possibilita aos estudantes ampliar sua compreensão sobre a vida, o nosso planeta e o universo, bem como sua capacidade de refletir, \textbf{argumentar}, propor soluções e enfrentar desafios pessoais e coletivos, locais e globais ( BRASIL, 2018c, \textbf{pp} 471, 472 ).
Nessa perspectiva, a Resolução n. 3 ( BRASIL, 2018a, artigo 11, §1º ), \textbf{que} atualiza as Diretrizes Curriculares para o Ensino Médio, dispõe sobre a organização do currículo por área de conhecimento, o \textbf{que} estimula o fortalecimento das relações entre os saberes e a sua \textbf{contextualização} para apreensão e intervenção na realidade, e requer planejamento e execução conjugados e colaborativos dos professores, com a valorização da \textbf{interdisciplinaridade} e a integração curricular na escola.
Dessa maneira, a área de CNT, na busca pela relação entre as competências específicas e habilidades, em continuidade à proposta do Ensino Fundamental, propõe \textbf{um} aprofundamento nas temáticas : Matéria e Energia/Vida, Terra e Cosmos/Processos e Práticas de Investigação.
Esses eixos temáticos permitem ao estudante investigar, analisar e discutir situações-problema relacionadas a diferentes contextos socioculturais, compreender leis, \textbf{teorias} e modelos para aplicá -los na resolução de problemas individuais, sociais e ambientais.
Cabe ressaltar \textbf{que} a ciência e a tecnologia estão fortemente inseridas no contexto \textbf{atual}, seja por meio dos seus impactos e das consequências presentes no cotidiano ou pelos produtos e serviços.
Alguns exemplos são os alimentos e medicamentos provenientes de organismos geneticamente modificados, transgênicos e produzidos em laboratórios, as tecnologias digitais, como cartões inteligentes, Home baking, palms, comércio eletrônico, voto eletrônico, dentre outros.
Essa explosão científica e tecnológica tem \textbf{exigido} cidadãos críticos e atuantes na resolução de problemas pertinentes à sua realidade.
Diante desse contexto, cabe às instituições educativas o papel de contribuir com o letramento científico e tecnológico dos estudantes e possibilitar o exercício pleno da sua cidadania.
Desse modo, ao considerar a relevância da “ Ciência, Tecnologia, Sociedade e Ambiente ( CTSA ) ” no contexto educacional e relacioná -la ao currículo, oportuniza -se o desenvolvimento de atividades e ações pedagógicas \textbf{que} permitam ao estudante perceber \textbf{que} a CTSA pode ser compreendida como produto construído em sociedade e \textbf{que} valoriza a construção da consciência cidadã, \textbf{baseada} no pensamento crítico, além de considerar os princípios éticos, os valores e a visão de mundo, para \textbf{que} esse \textbf{educando} possa participar democraticamente da proposição e tomada de decisões frente aos desafios Científicos, Sociais, Tecnológicos e Ambientais ( LEAL, 2009 ).
Nesse sentido, a BNCC propõe aos estudantes a utilização das Tecnologias Digitais de Informação e Comunicação ( TDIC ) de forma a relacionar os espaços colaborativos no ensino da Ciência com o uso de softwares na mediação do conteúdo, ou ainda, desfrutar de laboratórios virtuais, ambientes virtuais de aprendizagem, vídeos, redes sociais, fóruns, jogos digitais educativos, dentre outros.
A criação de atividades e a organização pedagógica precisam ter bases \textbf{metodológicas} \textbf{que} potencializam a aprendizagem, além de favorecer o desenvolvimento das competências socioemocionais de forma intencional, pois a concepção de educação integral deve considerar o estudante, em todas as suas dimensões e, nesse sentido, propor ações curriculares \textbf{que} respondam a essa proposição formativa.
Para \textbf{que} se alcance \textbf{um} caminho promissor, as metodologias diferenciadas, tais como : Metodologia de \textbf{Problematização}, Aprendizagens \textbf{baseadas} em Projetos, em Problemas, em Pares ou em Times, os Três Momentos Pedagógicos ( \textbf{Problematização} Inicial, Organização do Conhecimento, Aplicação do Conhecimento ) e Atividades Experimentais Investigativas possibilitam o desenvolvimento das competências gerais conforme propõe a BNCC.
Nesse contexto, em articulação com as competências gerais da Educação Básica no Currículo de Referência do Estado Mato Grosso do Sul - Ensino Médio, a área de CNT prima pela busca do letramento científico, \textbf{que} deve valorizar o conhecimento prévio dos estudantes, suas observações e constatações da realidade, próprios do saber empírico, de tal forma \textbf{que} proporcione condições para a (re)construção do conhecimento, \textbf{que} vai além de encontrar respostas científicas para solucionar problemas e explicar os fenômenos da natureza, mas de oferecer aos estudantes condições de compreender o sentido da vida e das leis \textbf{que} a regem.
Além disso, é \textbf{preciso} garantir aos estudantes o desenvolvimento de três competências específicas, organizadas em habilidades, a serem alcançadas nessa etapa.
Na competência específica 1, os fenômenos naturais e os processos tecnológicos são analisados sob a perspectiva das relações entre o eixo temático Matéria e Energia ; na competência específica 2 os processos de transformação e evolução \textbf{permeiam} a natureza e ocorrem das moléculas às estrelas em diferentes escalas de tempo.
Por meio dessas competências, os estudantes têm a oportunidade de elaborar reflexões \textbf{que} situem a \textbf{humanidade} e o planeta Terra na história do Universo por meio do eixo temático Vida, Terra e Cosmos.
Por fim, a competência específica 3 articula em suas habilidades a proposta de \textbf{um} mundo repleto de informações de diferentes naturezas e origens, \textbf{que} são facilmente difundidas e acessadas, sobretudo, por meios digitais.
Assim, é premente \textbf{que} os estudantes desenvolvam capacidades de seleção e discernimentos de informações \textbf{que} lhes permitam, com base em conhecimentos científicos confiáveis, investigar situações-problema e avaliar as aplicações desses conhecimentos nas diversas esferas da vida humana, com ética e responsabilidade, por meio do eixo temático Processos e Práticas de Investigação.
Conforme consta na BNCC : 1.
Analisar fenômenos naturais e processos tecnológicos, com base nas interações e relações entre matéria e energia, para propor ações individuais e coletivas \textbf{que} aperfeiçoem processos produtivos, minimizem impactos socioambientais e melhorem as condições de vida em âmbito local, regional e global. 2.
Analisar e utilizar interpretações sobre a dinâmica da Vida, da Terra e do Cosmos para elaborar argumentos, realizar previsões sobre o funcionamento e a evolução dos seres vivos e do Universo, e fundamentar e defender decisões éticas e responsáveis. 3.
Investigar situações-problema e avaliar aplicações do conhecimento científico e tecnológico e suas implicações no mundo, utilizando procedimentos e linguagens próprios das Ciências da Natureza, para propor soluções \textbf{que} considerem demandas locais, regionais e/ou globais, e comunicar suas descobertas e conclusões a públicos variados, em diversos contextos e por meio de diferentes mídias e tecnologias digitais de informação e comunicação ( TDIC ). ” ( BRASIL, 2018c, p. 553 ) A área de CNT, composta pela Biologia, Física e Química, está integrada do ponto de vista científico e tecnológico, \textbf{uma} vez \textbf{que} é possível estabelecer relações entre os conceitos e contextos de aplicação teórico-prática para o desenvolvimento e avanço da cultura científica.
Nessa concepção, destacam -se algumas especificações em cada componente curricular a partir das possibilidades \textbf{metodológicas} e didáticas para intervenção pedagógica \textbf{que} proporcione a formação humana integral e o protagonismo do estudante.
O componente curricular Biologia é responsável por estudar a vida, destacando seu surgimento, processo evolutivo, constituição e as formas de interação entre os fatores bióticos e abióticos.
A Biologia viabiliza o percurso por vários temas, tais como saúde, genética, meio ambiente, sistemática, taxonomia, evolução dos seres vivos e outras diversas possibilidades, \textbf{que} facilitam a compreensão das relações e a \textbf{interdisciplinaridade} com outras áreas de conhecimento, como Ciências Humanas e Sociais Aplicadas, Matemática e Linguagens, possibilitando a \textbf{influência} e a interação em todas as formas de vida e na organização da sociedade.
Em Física o estudo de diversos acontecimentos do dia a dia, como os fenômenos naturais e os processos tecnológicos, proporciona aos estudantes compreender a natureza com ênfase na admiração pela ciência nas suas diversas áreas, como a comparação entre Mecânica Clássica e Quântica, os efeitos das Ondas Sonoras e Eletromagnéticas na vida das pessoas, assim como a importância de compreender de forma estruturada o conhecimento formal da física por meio da leitura de artigos científicos.
No Ensino Médio esse componente propõe \textbf{uma} aprendizagem com compreensão e capacidade de aplicação, com \textbf{métodos} \textbf{que} incluem a investigação, \textbf{contextualização} e \textbf{interdisciplinaridade}.
Proporciona \textbf{uma} formação \textbf{que} engloba os temas \textbf{atuais} do mundo \textbf{contemporâneo} e ações \textbf{que} permitam aos estudantes analisar variáveis, coletar dados e formular explicações com representações, simulações e construção de protótipos, acompanhadas de diálogos constantes entre os envolvidos.
A ciência, \textbf{que} se dedica ao estudo e a análise da matéria, sua composição, transformações e energia envolvidas nos processos, denomina -se Química e busca, do ponto de vista microscópico ( partículas, átomos e moléculas ) e macroscópico ( substâncias e materiais ), entender os fenômenos, sejam eles químico ou físicos, a proposição de modelos explicativos \textbf{que} possam ampliar e consolidar a \textbf{percepção} da realidade em \textbf{que} se \textbf{vive}.
Para o desenvolvimento do componente curricular Química é importante a adequação da linguagem por meio de temáticas \textbf{que} relacionam a sociedade e o cotidiano do estudante, tais como : Biocombustíveis, Saneamento Básico, Catalisadores, Farmoquímicos, Petroquímica, Química Verde, dentre outras, e \textbf{que} as ações pedagógicas considerem a reflexão epistemológica na aprendizagem, promovendo \textbf{um} ambiente educacional \textbf{que} correlacione a \textbf{teoria} e a prática dos conceitos químicos e a compreensão dos fenômenos materiais, de modo \textbf{que} esse \textbf{educando} conheça e interprete os aspectos do mundo \textbf{que} afetam sua vida diária ( LEAL, 2009 ).
Além disso, a Química possibilita ao estudante a capacidade de manifestar -se sobre os desafios cotidianos para \textbf{que} possa \textbf{aprimorar} as atividades básicas para a melhoria da vida humana, de forma a subsidiar a formação humana integral por meio de \textbf{uma} cultura e prática científica.
Nesse contexto, a área Ciências da Natureza pretende formar cidadãos mais críticos, conscientes de seu papel social, político, econômico, socioambiental e ambiental, facilitando o acesso às novas tecnologias e às descobertas científicas, de forma contextualizada, dando ao objeto de conhecimento estudado \textbf{uma} \textbf{aplicabilidade} para a vida.
Deve -se desenvolver \textbf{uma} prática reflexiva \textbf{que} amplie as possibilidades de aprendizagem, para \textbf{que} se encontre sentido no \textbf{que} se apreende.
Além disso, as aprendizagens dos componentes curriculares dessa área de conhecimento devem estar articuladas com o \textbf{que} cerca o estudante, valorizando seus conhecimentos prévios e a sua visão do mundo, para \textbf{que} possa decodificar essas informações, tais como : as questões emergentes do aquecimento global, o uso de tecnologias, os avanços da medicina, as discussões sobre células-tronco embrionárias e o seu papel na cidadania.
Essas discussões estão no campo do letramento científico, pois as informações e as propostas de pesquisas possibilitam \textbf{que} o estudante ressignifique a sua leitura de mundo.
Assim, ao desenvolver as atividades pedagógicas do ponto de vista curricular, a \textbf{teoria} e a prática devem promover a aprendizagem de conceitos, práticas científicas e a pesquisa na busca por resoluções de questionamentos e problemas, sejam eles de caráter político, social, econômico, socioambiental ou ambiental, considerando a historicidade dos fenômenos na perspectiva de combater a \textbf{fragmentação} do conhecimento e oportunizar o protagonismo em ações \textbf{que} possam contribuir para a aplicação de novas tecnologias na área das CNT, fundamentais para a sociedade.
Por meio dessa \textbf{contextualização}, apresenta -se o organizador curricular da área de CNT, destacando -se : as Competências Específicas da área, \textbf{que} norteiam os Eixos Temáticos e organizam as habilidades para mobilizar as aprendizagens e integrar os componentes curriculares e seus objetos de conhecimento ; e as Sugestões Didáticas, \textbf{que} servem como aporte pedagógico para os professores na compreensão do Currículo de Referência do Estado do Mato Grosso do Sul - Ensino Médio e na sua operacionalização por meio de práticas \textbf{que} garantam de forma intencional a aprendizagem e formação humana e integral do estudante.

% matched lemmas: agregar, aplicabilidade, aprimorar, argumentar, atual, basear, contemporâneo, contextualização, desdobramento, despertar, educar, exigir, fragmentação, humanidade, influência, inter-relação, interdisciplinaridade, metodológico, método, percepção, permear, pp, preciso, problematização, que, século, teoria, teórico, um, viver
\end{textsample}
